\chapter*{Abstract}
\addcontentsline{toc}{chapter}{Abstract}

SavePoint grew out of applying to video games the personal cataloguing pattern of platforms such as Goodreads or Letterboxd, adding an inventory of physical and digital ownership that those products never had to solve; on that basis, the project turns the platform into a reproducible testbed for recommender systems. The platform includes a catalogue, library, physical and digital copies, ratings, lists, comments, profiles, privacy, social relations and recommendations. The application and the reproducible evaluation of the recommenders are the two complementary halves of the work: the former collects and governs the data the latter needs.

The system compares sixteen reference algorithms: baselines, content-based variants, neighbourhood collaborative filtering, and hybrid approaches with diversity re-ranking. This report describes item and user representations, tag weighting, facet similarity, external quality, popularity, negative evidence, cold start, and hybrid combination. The offline protocol fixes corpus, snapshots, seeds, shared candidates, exclusions and ranking metrics at $K\in\{5,10,20\}$.

The published v15 evidence uses 400 active synthetic users and 79 evaluable users, with a single test run. It supports an internal comparison under the frozen corpus and protocol, but it does not represent real users or justify automatic generalisation. This report explicitly separates that publication from the later implementation contract.

An agent-assisted engineering methodology governed through the Get Stuff Done method and conceptually organised with Obsidian was applied throughout the project. The tools were used under human direction and review, while final decisions and responsibility for analysis, design, implementation and validation remained with Felipe Peña Núñez.

\vspace{.5cm}
\textbf{Keywords:} video games, recommender systems, offline evaluation, reproducibility, collaborative filtering, content-based recommendation, data provenance.
